\documentclass[11pt]{article}
\usepackage[margin=1in]{geometry}
\usepackage{amsmath,amssymb,amsthm}
\usepackage{array}
\usepackage{parskip}
\usepackage{booktabs}

\newtheorem{theorem}{Theorem}
\newtheorem{lemma}[theorem]{Lemma}
\newtheorem{corollary}[theorem]{Corollary}
\newtheorem{proposition}[theorem]{Proposition}
\theoremstyle{definition}
\newtheorem{definition}[theorem]{Definition}
\theoremstyle{remark}
\newtheorem{remark}[theorem]{Remark}

\setlength{\emergencystretch}{3em}

\newcommand{\F}{\mathbb{F}}
\newcommand{\ceil}[1]{\left\lceil #1 \right\rceil}

\title{Disproof of conjecture \texttt{00000001067}}
\author{earthking11}
\date{2026-09-14}

\begin{document}
\maketitle

\section*{Verdict: FALSE}

We disprove conjecture \texttt{00000001067} as stated. The refutation is
unconditional, completely elementary, and finitely verifiable: for the prime
$p = 19$, which satisfies $p \equiv 3 \pmod 4$, the largest strong Sidon set in
$\F_{19}$ has size $4$, whereas $\ceil{\sqrt{19}} = 5$. The filed clause that
the $O(1)$ term equals $0$ for $p \equiv 3 \pmod 4$ is an \emph{exact} claim,
and it is false. It fails already at $p = 11$ and also at $p = 43, 47, 59$.

\section{The conjecture, quoted verbatim}

From \texttt{conjectures/00000001067.md}:

\begin{quote}
\textbf{English.} Definition: A $B_h$ set (unique $h$-fold sums). Conjecture:
The maximal $B_2$ set in $\F_p$ has size $\ceil{\sqrt p} + O(1)$, and the
$O(1)$ term equals $0$ for $p \equiv 3 \bmod 4$ (an exact $B_2$ result).
\end{quote}

The conjecture is filed in both English and Chinese; the Chinese original is
reproduced verbatim in Unicode in \texttt{README.md}, and the two versions are
identical in content. The PDF font does not embed CJK glyphs, so the Chinese is
kept in the markdown copy rather than duplicated here.

The phrase ``an exact $B_2$ result'' in the last clause fixes the reading: the
$O(1)$ term is asserted to \emph{equal} $0$, so the conjecture claims the exact
identity
\begin{equation}\label{eq:claim}
\max\{\,|S| : S \subseteq \F_p \text{ is a } B_2 \text{ set}\,\}
   = \ceil{\sqrt p}
\qquad \text{for every prime } p \equiv 3 \pmod 4 .
\end{equation}
We refute \eqref{eq:claim}.

\section{Convention: $B_2$ means strong Sidon}

\begin{definition}[Strong Sidon / $B_2$ set]\label{def:strong}
A subset $S \subseteq \F_p$ is a \emph{(strong) Sidon set}, equivalently a
$B_2$ set with unique $2$-fold sums, if the sums $a+b$ with $a,b \in S$ and
$a \le b$ (repetition allowed) are pairwise distinct. Equivalently, the map
$\{a,b\} \mapsto a+b$ from unordered pairs with repetition to $\F_p$ is
injective.
\end{definition}

This is the standard meaning of ``unique $2$-fold sums'' and of $B_2$: a
$B_h$ set has all $h$-fold sums (unordered, repetitions allowed) distinct, so a
$B_2$ set is exactly a strong Sidon set as in
Definition~\ref{def:strong}. In particular $2a = b + c$ with $b \ne c$ is
forbidden, which is what distinguishes this from the \emph{weak} Sidon reading
discussed in Section~\ref{sec:readings}.

\section{The difference-count proof}

\begin{lemma}[Difference bound]\label{lem:diff}
Let $S \subseteq \F_p$ be a strong Sidon set of size $k$. Then the $k(k-1)$
ordered differences $a - b$ with $a,b \in S$, $a \ne b$, are pairwise distinct
and nonzero. Consequently
\[
k(k-1) \;\le\; p - 1 .
\]
\end{lemma}

\begin{proof}
Let $a \ne b$ and $c \ne d$, and suppose $a - b = c - d$. Then
$a + d = c + b$. Both sides are sums of two elements of $S$ with repetition
allowed. If the multisets $\{a,d\}$ and $\{c,b\}$ are equal, then
$(a,d) = (c,b)$: indeed $\{a,d\} = \{c,b\}$ means either $a=c$, $d=b$, or
$a=b$, $d=c$, and the second alternative is excluded by $a \ne b$. In the first
case $a-b = a-b$, i.e.\ $(a,b) = (c,d)$. If the two multisets differ, then
$a+d = c+b$ is a collision between two distinct allowed sums, contradicting
Definition~\ref{def:strong}. Hence the differences are pairwise distinct.
They are nonzero because $a \ne b$. Finally, $\F_p$ has exactly $p-1$ nonzero
elements, so $k(k-1) \le p-1$.
\end{proof}

\begin{corollary}\label{cor:19}
There is no strong Sidon set of size $5$ in $\F_{19}$, so the maximum is at
most $4$. More generally, at $p = 19$ one has $5 \cdot 4 = 20 > 18 = p-1$.
\end{corollary}

\begin{proof}
If $|S| = 5$, Lemma~\ref{lem:diff} would give $5\cdot 4 = 20 \le 19 - 1 = 18$,
a contradiction.
\end{proof}

\begin{proposition}\label{prop:max4}
The maximum size of a strong Sidon set in $\F_{19}$ is exactly $4$.
\end{proposition}

\begin{proof}
The set $\{0,1,3,7\}$ is strong Sidon: its ten unordered sums are
\[
0,\;1,\;3,\;7,\;\;2,\;4,\;8,\;\;6,\;10,\;\;14
\]
(reading $\{0+0, 0+1, 0+3, 0+7, 1+1, 1+3, 1+7, 3+3, 3+7, 7+7\}$), i.e.\ the
sorted list
\[
0,\,1,\,2,\,3,\,4,\,6,\,7,\,8,\,10,\,14 ,
\]
which is pairwise distinct. Hence the maximum is at least $4$;
Corollary~\ref{cor:19} gives at most $4$. Therefore it equals $4$.
\end{proof}

\begin{theorem}\label{thm:false}
Conjecture \texttt{00000001067} is false.
\end{theorem}

\begin{proof}
At the prime $p = 19$ we have $p \equiv 3 \pmod 4$ (indeed $19 = 4\cdot 4 + 3$)
and $\ceil{\sqrt{19}} = 5$, because $4^2 = 16 < 19 \le 25 = 5^2$. By
Proposition~\ref{prop:max4} the maximal strong Sidon set in $\F_{19}$ has size
$4$, not $5$. Thus the exact identity \eqref{eq:claim} fails at $p = 19$, which
refutes the ``$O(1)$ term equals $0$ for $p \equiv 3 \bmod 4$'' clause.
\end{proof}

\begin{remark}
$p = 11$ is an even smaller counterexample: $11 \equiv 3 \pmod 4$,
$\ceil{\sqrt{11}} = 4$, and Lemma~\ref{lem:diff} forbids $k = 4$ because
$4 \cdot 3 = 12 > 10 = 11 - 1$; a maximum $3$-set is $\{0,1,3\}$. We use
$p = 19$ as the headline witness because the difference bound there is
$5 \cdot 4 = 20 > 18 = p-1$.
\end{remark}

\section{Exhaustive verification at $p = 19$}

The counting argument above is decisive, but the claim can also be checked
directly. Enumerating all $\binom{19}{5} = 11628$ five-element subsets of
$\F_{19}$ and testing the $15$ unordered sums of each, one finds that
\emph{every} five-subset contains a repeated sum: there is no strong Sidon
$5$-set, confirming Corollary~\ref{cor:19} without the difference bound.
\texttt{reproduce.py} performs exactly this enumeration. It also reports that
all $2129$ strong Sidon subsets of $\F_{19}$ have size at most $4$, and that
the maximum is attained (a maximum $4$-set is $\{0,1,3,7\}$).

\section{The full table}

We computed the maximum strong Sidon size in $\F_p$ by exhaustive backtracking
for $p = 7, 11, 19, 23,$ $31, 43, 47, 59$. Column ``strong'' is the true
maximum; ``$\ceil{\sqrt p}$'' is the conjectured exact value.

\begin{center}
\begin{tabular}{r r r r l}
\toprule
$p$ & $p \bmod 4$ & $\ceil{\sqrt p}$ & max strong Sidon & holds? \\
\midrule
$7$  & $3$ & $3$ & $3$ & yes \\
$11$ & $3$ & $4$ & $3$ & \textbf{NO} \\
$19$ & $3$ & $5$ & $4$ & \textbf{NO} \\
$23$ & $3$ & $5$ & $5$ & yes \\
$31$ & $3$ & $6$ & $6$ & yes \\
$43$ & $3$ & $7$ & $6$ & \textbf{NO} \\
$47$ & $3$ & $7$ & $6$ & \textbf{NO} \\
$59$ & $3$ & $8$ & $7$ & \textbf{NO} \\
\bottomrule
\end{tabular}
\end{center}

The exact formula \eqref{eq:claim} happens to hold at $p = 7, 23, 31$ but fails
at $p = 11, 19, 43, 47, 59$. This is what makes the conjecture plausible but
false rather than obviously wrong: the values agree on the first few primes
$p \equiv 3 \pmod 4$ that one is likely to test, and only break later. For
$p = 43, 47, 59$ the difference bound alone does not forbid the conjectured
size: at $p = 43$ one has $7 \cdot 6 = 42 = 43 - 1$, at $p = 47$,
$7 \cdot 6 = 42 \le 46$, and at $p = 59$, $8 \cdot 7 = 56 \le 58$; it is the
exhaustive search (and, for $p=43$, the Bruck--Ryser--Chowla theorem below)
that rules the larger size out.

\section{Cross-check at $p = 43$: Bruck--Ryser--Chowla}

At $p = 43$ a strong Sidon set of size $k = 7$ would satisfy
$k(k-1) = 42 = p - 1$, so by Lemma~\ref{lem:diff} its $42$ differences would be
\emph{all} the nonzero elements of $\F_{43}$; such a set is a perfect
$(43,7,1)$ difference set. A $(v,k,\lambda)$ difference set with $k-\lambda=6$
gives a symmetric $2$-$(v,k,\lambda)$ design, and a symmetric $2$-$(n^2+n+1,
n+1, 1)$ design is a projective plane of order $n$. Here
$v = 43 = 6^2 + 6 + 1$ and $k = 7 = 6 + 1$, so a size-$7$ strong Sidon set in
$\F_{43}$ would produce a projective plane of order $6$.

The Bruck--Ryser--Chowla theorem states that if a projective plane of order
$n$ exists and $n \equiv 1, 2 \pmod 4$, then $n$ is a sum of two squares. Here
$6 \equiv 2 \pmod 4$, and $6$ is not a sum of two squares ($0,1,2,4,5$ are
representable below $6$, but $6$ is not). Hence no projective plane of order
$6$ exists, so no size-$7$ strong Sidon set in $\F_{43}$ exists. This is
consistent with the computed maximum $6$, and it independently certifies the
failure of \eqref{eq:claim} at $p = 43$.

\section{Alternative readings}\label{sec:readings}

We tried to save the conjecture by changing the reading; none succeeds.

\subsection*{Weak Sidon}

\begin{definition}[Weak Sidon]
$S \subseteq \F_p$ is \emph{weak Sidon} if the sums $a+b$ with $a,b \in S$,
$a \ne b$ (distinct elements only) are pairwise distinct.
\end{definition}

This is genuinely weaker: repetitions $2a$ are unconstrained. The weak maxima
for the same primes are

\begin{center}
\begin{tabular}{r r r r}
\toprule
$p$ & $\ceil{\sqrt p}$ & max weak Sidon & holds? \\
\midrule
$7$  & $3$ & $4$  & \textbf{NO} \\
$11$ & $4$ & $5$  & \textbf{NO} \\
$19$ & $5$ & $6$  & \textbf{NO} \\
$23$ & $5$ & $6$  & \textbf{NO} \\
$31$ & $6$ & $7$  & \textbf{NO} \\
$43$ & $7$ & $8$  & \textbf{NO} \\
$47$ & $7$ & $8$  & \textbf{NO} \\
$59$ & $8$ & $9$  & \textbf{NO} \\
\bottomrule
\end{tabular}
\end{center}

Under the weak reading the exact formula fails even more dramatically: at
$p = 19$ the maximum is $6 > 5 = \ceil{\sqrt{19}}$ (for instance
$\{0,1,2,4,7,12\}$ works), and at $p = 11$ it is $5 > 4 = \ceil{\sqrt{11}}$
(for instance $\{0,1,2,4,7\}$). Note also that Lemma~\ref{lem:diff} is a
\emph{strong}-Sidon statement: the weak maximum set $\{0,1,2,4,7,12\}$ in
$\F_{19}$ has $6 \cdot 5 = 30$ ordered differences but $\F_{19}$ has only $18$
nonzero elements, so its differences are certainly not distinct. The
difference-count proof must not be applied to weak Sidon sets; for weak Sidon
the correct counting bound is $\binom{k}{2} \le p$, since the $\binom{k}{2}$
sums of distinct pairs must lie in $\F_p$.

\subsection*{``Maximal'' as inclusion-maximal}

If ``maximal'' is read as \emph{inclusion-maximal} (a strong Sidon set that is
not properly contained in another strong Sidon set) rather than maximum
cardinality, nothing changes: every inclusion-maximal strong Sidon set in
$\F_{19}$ has size $4$. There are $1140$ such sets, and all have size $4$;
none has size $5$. So this reading does not rescue the $O(1) = 0$ clause
either.

\subsection*{Asymptotic weakening}

Only a substantially weakened asymptotic version survives. The classical
upper bound for strong Sidon sets in $\F_p$ is $|S| \le \sqrt p + O(1)$
(immediately from $k(k-1) \le p-1$), so the asymptotic part of the conjecture
is true and standard; what fails is precisely the exact ``$O(1)$ term $= 0$''
clause, i.e.\ the assertion that the maximum equals $\ceil{\sqrt p}$ for
$p \equiv 3 \pmod 4$. Replacing the exact claim by $|S| = \ceil{\sqrt p} +
O(1)$ would make it true but vacuous as a statement about the $O(1)$ term.

\section{Reproducibility}

The exhaustive computations are carried out by \texttt{reproduce.py} using the
Python~3 standard library only:

\begin{quote}\small\ttfamily
\$ python3 reproduce.py\\
... \textit{(checks, the strong/weak tables,}\\
\hspace*{1.5em}\textit{difference bound,}\\
\hspace*{1.5em}\textit{Bruck--Ryser cross-check)}\\
PASS: all checks verified;\\
\hspace*{1.5em}conjecture 00000001067 is FALSE.
\end{quote}

It computes the strong and weak maxima for all eight primes by exhaustive
backtracking, verifies the difference bound on every maximum set, verifies the
witness $\{0,1,3,7\}$ by listing its sums, enumerates all $\binom{19}{5} =
11628$ five-subsets, checks the equivalence ``strong Sidon $\Leftrightarrow$
differences distinct'' on all subsets of $\F_7, \F_{11}$ and all subsets of
size $\le 6$ of $\F_{19}, \F_{23}$, and reports the Bruck--Ryser--Chowla and
inclusion-maximal facts. It prints \texttt{PASS} and exits $0$ exactly when
every check holds.

\section{Formalisation}

The refutation is formalised in the Lean~4 project under \texttt{lean4/}
(core Lean only, \texttt{import Std}, no Mathlib, no \texttt{sorry}). Over the
field $\F_{19}$ presented as \texttt{Fin 19}, the main statements are

\begin{quote}\small\ttfamily
theorem no\_sidon5 :\\
\hspace*{2em}\(\forall\) a b c d e : Fin 19,\\
\hspace*{2em}\phantom{\(\forall\) }a < b \(\to\) b < c \(\to\) c < d
   \(\to\) d < e\\
\hspace*{2em}\phantom{\(\forall\) }\(\to\) sidon5 a b c d e = false\\
theorem sidon4\_witness :\\
\hspace*{2em}sidon4 0 1 3 7 = true\\
theorem count\_obstruction :\\
\hspace*{2em}5 * 4 > 19 - 1\\
theorem ceil\_sqrt19 :\\
\hspace*{2em}4 \^{} 2 < 19 \(\wedge\) 19 \(\le\) 5 \^{} 2\\
theorem conjecture\_00000001067\_false :\\
\hspace*{2em}\((\exists\) a b c d : Fin 19,\\
\hspace*{2em}\phantom{(\(\exists\) }sidon4 a b c d = true)\\
\hspace*{2em}\(\wedge\) no\_sidon5 \(\wedge\) ceil\_sqrt19
\end{quote}

Here \texttt{sidon5 a b c d e} decides whether the $15$ unordered sums
$a+a,\dots,e+e$ are pairwise distinct (strong Sidon), and \texttt{sidon4} is
the analogous predicate for four elements. The proof of \texttt{no\_sidon5}
uses translation invariance of the strong Sidon property: replacing each
element $x$ by $x-a$ shifts every sum by $-2a$, an injection on $\F_{19}$,
so it suffices to refute the case $a=0$, which \texttt{decide} checks
exhaustively; the general case is then reduced to it. The audit
\texttt{Check.lean} reports no \texttt{sorryAx} and no \texttt{Lean.ofReduceBool}
(the latter would signal \texttt{native\_decide}, which is deliberately not
used); the only axioms are the core \texttt{propext} and \texttt{Quot.sound}.
See \texttt{lean4/README.md} for the full statement table and the honest axiom
footprint.

\section*{Status against the submission rules}

Rule 3 requires a LaTeX source, a PDF document, and a Lean~4 project.

\begin{itemize}
\item \textbf{LaTeX source} --- \texttt{main.tex}, a standalone \texttt{article}
  using \texttt{geometry}, \texttt{amsmath}, \texttt{amssymb}, \texttt{amsthm},
  \texttt{array}, \texttt{parskip}, \texttt{booktabs}; compiles with
  \texttt{tectonic main.tex}.
\item \textbf{PDF} --- at \texttt{build/main.pdf}, produced by
  \texttt{tectonic --outdir build main.tex}.
\item \textbf{Lean~4 project} --- \texttt{lean4/} with \texttt{lean-toolchain}
  pinned to \texttt{leanprover/lean4:v4.33.1}, \texttt{lakefile.toml}
  (library \texttt{Main}, no dependencies), \texttt{Main.lean},
  \texttt{Check.lean} (\texttt{\#print axioms}), and \texttt{lean4/README.md}.
  \texttt{lake build} exits $0$ and the audit reports no \texttt{sorryAx}.
\end{itemize}

\end{document}
